% ============ 第6章 §6.2 Littlewood--Paley 理论的应用 (原书 pp.130-132) ============

\section{Littlewood--Paley 理论的应用}\label{sec:6.2}
\index{R@Riesz 变换}




这里, 设 $m$ 是在除去原点的 $\mathbb{R}^{n}$ 上定义的 $C^{\infty}$ 级函数, 对所有多重指数 $\alpha$ 满足 $|\partial^{\alpha}m(\xi)| \lesssim |\xi|^{-|\alpha|}$. 不过, $\lesssim$ 中所含的常数可以依赖于 $\alpha$. $T_{m}f = \mathcal{F}^{-1}[m\mathcal{F}f]$ 由 Plancherel 定理是 $L^{2}$ 有界的, 但要对 $1 < p < \infty$ 证明 $T_{m}$ 可以扩张为 $L^{p}$ 有界线性算子.

\begin{theorem}\label{thm:6.2}
对 $f \in \mathcal{S}_{\infty}(\mathbb{R}^{n})$, $T_{m}f \in \mathcal{S}_{\infty}(\mathbb{R}^{n})$ 成立, 且若 $1 < p < \infty$, 则成立不依赖于 $f \in \mathcal{S}_{\infty}(\mathbb{R}^{n})$ 的常数估计 $\|T_{m}f\|_{L^{p}} \lesssim \|f\|_{L^{p}}$.
\end{theorem}

作为这样的 $m$ 的例子, 对 $j = 1, 2, \ldots, n$ 有
\[
m(\xi) = \frac{\xi_{j}}{|\xi|},
\]
与这个 $m$ 对应的 $T_{m}$ 称为第 $j$ Riesz 变换.

\begin{proof}
设 $g \in \mathcal{S}(\mathbb{R}^{n})$. $g \in \mathcal{S}_{\infty}(\mathbb{R}^{n})$ 与 $\partial^{\alpha}\mathcal{F}g(0) = 0$, $\alpha \in \mathbb{N}_{0}^{n}$ 是等价的. 因此, 使用这个判定, 可得 $T_{m}f \in \mathcal{S}_{\infty}(\mathbb{R}^{n})$.

来证明范数不等式 $\|T_{m}f\|_{L^{p}} \lesssim \|f\|_{L^{p}}$. 为此, 只要证明
\[
\left\|\left(\sum_{j=-\infty}^{\infty}|\varphi_{j}(D)T_{m}f|^{2}\right)^{1/2}\right\|_{L^{p}} \lesssim \|f\|_{L^{p}}
\]
即可. 通过使用 Rademacher 序列来考虑, 只要证明
\[
L_{t}(z) = \sum_{j=-J}^{J}r_{j}(t)\mathcal{F}^{-1}[m\cdot\varphi_{j}](z)
\]
满足以下条件即可.

(1) $\|\mathcal{F}L_{t}\|_{L^{\infty}} \lesssim 1$.

(2) $\sup_{z \in \mathbb{R}^{n}, t \in (0,1)}|z|^{n}|L_{t}(z)| \lesssim 1$.

(3) $\sup_{z \in \mathbb{R}^{n}, t \in (0,1)}|z|^{n+1}|\nabla L_{t}(z)| \lesssim 1$.

这与证明 Littlewood--Paley 理论时所考虑的 $K_{t}$ 的估计相同, 因此省略细节.
\end{proof}

\begin{exercise}\label{exer:6.2}
若 $m \in L^{\infty}(\mathbb{R}^{n})$, 则证明 $T_{m}$ 是 $L^{2}(\mathbb{R}^{n})$ 有界的.
\end{exercise}

\begin{exercise}\label{exer:6.3}
设 $m$ 是在除去原点的 $\mathbb{R}^{n}$ 上定义的 $C^{\infty}$ 级函数, 对所有多重指数 $\alpha$ 满足 $|\partial^{\alpha}m(\xi)| \lesssim |\xi|^{-|\alpha|}$. 对 $g \in \mathcal{S}_{\infty}(\mathbb{R}^{n})$, 证明 $T_{m}g \in \mathcal{S}_{\infty}$.
\end{exercise}
